\setcounter{aufgabe}{0}
\hypertarget{e1}{}\einheitenkopf{Einheit 1 von 5 · Wurzelziehen und Lösbarkeit}

\begin{aufgabe}{Wurzelziehen -- Quadratisch oder linear? Kreuze an, ob in der Gleichung
ein $x^2$ oder eine Klammer mit $x$ im Quadrat vorkommt. Löse nichts.}
\begin{teile}
\teil $4x - 9 = 15$ \hfill \kreuz{quadratisch}\kreuz{linear}
\teil $x^2 = 121$ \hfill \kreuz{quadratisch}\kreuz{linear}
\teil $(x - 5)^2 = 9$ \hfill \kreuz{quadratisch}\kreuz{linear}
\teil $2x + x = 33$ \hfill \kreuz{quadratisch}\kreuz{linear}
\teil $x \cdot (x + 7) = 0$ \hfill \kreuz{quadratisch}\kreuz{linear}
\end{teile}
\end{aufgabe}

\begin{aufgabe}{Wurzelziehen -- Zwei, eine oder keine? Sieh nur auf die Zahl rechts und
kreuze an, wie viele Lösungen die Gleichung hat. Rechne nichts.}
\begin{teile}
\teil $x^2 = 169$ \hfill \kreuz{zwei}\kreuz{eine}\kreuz{keine}
\teil $x^2 = -12$ \hfill \kreuz{zwei}\kreuz{eine}\kreuz{keine}
\teil $x^2 = 0$ \hfill \kreuz{zwei}\kreuz{eine}\kreuz{keine}
\teil $x^2 = 5$ \hfill \kreuz{zwei}\kreuz{eine}\kreuz{keine}
\teil $x^2 = -0{,}5$ \hfill \kreuz{zwei}\kreuz{eine}\kreuz{keine}
\end{teile}
\end{aufgabe}

\begin{aufgabe}{Wurzelziehen -- rechnen. Zieh die Wurzel und gib beide Lösungen an.}
\beispiel{x^2 &= 49 &&\mid \text{Wurzel, beide Vorzeichen} \\
x_1 &= 7 && x_2 = -7}
\begin{gleichungsraster}[2]
\gl{x^2 = 9} & \gl{x^2 = 64} \\
\gl{x^2 = 121} & \gl{x^2 = 225} \\
\end{gleichungsraster}
Geht die Wurzel nicht auf, lass sie stehen und gib zusätzlich den Näherungswert auf
zwei Stellen an. Gibt es keine Lösung, schreib das hin.
\begin{gleichungsraster}[2]
\gl{x^2 = 11} & \gl{x^2 = -100} \\
\end{gleichungsraster}
Stell erst $x^2$ frei, dann zieh die Wurzel.
\begin{gleichungsraster}[3]
\gl{x^2 - 45 = 55} & \gl{3x^2 = 48} \\
\gl{2x^2 + 5 = 55} & \\
\end{gleichungsraster}
Rechne die Klammer rückwärts: erst die Wurzel mit beiden Vorzeichen, dann die Zahl
hinüberbringen.
\begin{gleichungsraster}[3]
\gl{(x - 5)^2 = 36} & \gl{(x + 3)^2 = 81} \\
\gl{(x - 7)^2 = 0} & \\
\end{gleichungsraster}
\end{aufgabe}

\begin{aufgabe}{Wurzelziehen -- Probe. Setze die Zahl für $x$ ein, rechne beide Seiten
aus und kreuze an.}
\begin{teile}
\teil $x^2 = 64$ \quad für $x = -8$ \hfill \kreuz{wahr}\kreuz{falsch}
\teil $(x - 5)^2 = 16$ \quad für $x = 1$ \hfill \kreuz{wahr}\kreuz{falsch}
\teil $(x + 3)^2 = 16$ \quad für $x = 2$ \hfill \kreuz{wahr}\kreuz{falsch}
\end{teile}
\end{aufgabe}

\begin{aufgabe}{Wurzelziehen -- ablesen. Die Parabel gehört zu $p(x) = x^2 - 3$. Lies
an der Zeichnung ab, wie viele Lösungen die Gleichung hat, und gib die Lösungen an.}

\begin{ksys}[xmin=-5,xmax=5,ymin=-6,ymax=7,ablesen]
\parabel{1}{0}{-3}{p}
\gerade[-4]{0}{1}{y=1}
\gerade[4]{0}{-3}{y=-3}
\gerade[-4]{0}{-5}{y=-5}
\end{ksys}

\begin{teile}
\teil $x^2 - 3 = 1$ \hfill \leerfeld Lösungen: \leerfeld
\teil $x^2 - 3 = -3$ \hfill \leerfeld Lösungen: \leerfeld
\teil $x^2 - 3 = -5$ \hfill \leerfeld Lösungen: \leerfeld
\end{teile}
\end{aufgabe}

\begin{aufgabe}{Wurzelziehen -- Fehler finden. Jonas löst so:}
\rechnung{(x + 6)^2 &= 25 \\ x_1 &= 6 + 5 = 11 \\ x_2 &= 6 - 5 = 1}
Finde den Fehler, benenne ihn und korrigiere die Rechnung.

\begin{teile}
\teil Löse $(x + 4)^2 = 81$.
\end{teile}
\end{aufgabe}

\begin{aufgabe}{Wurzelziehen -- Begründen. Antworte ohne zu rechnen.}
\begin{teile}
\teil Begründe, warum $x^2 = -20$ keine Lösung hat.
\teil Max sagt: „$x^2 = 0$ hat keine Lösung.`` Hat er recht? Begründe.
\teil Begründe, warum $x^2 = 7$ zwei Lösungen hat, obwohl $7$ keine Quadratzahl ist.
\end{teile}
\end{aufgabe}

\begin{aufgabe}{Wurzelziehen -- aus dem Sachtext. Rechne und antworte mit Einheit.}
\begin{teile}
\teil Ein quadratisches Beet hat einen Flächeninhalt von $64\,\text{m}^2$. Berechne
seine Seitenlänge. \hfill (P10 2020)
\teil Ein quadratisches Grundstück soll $200\,\text{m}^2$ groß werden. Passt es auf
eine freie Fläche von $14\,\text{m}$ mal $14\,\text{m}$? Rechne und antworte.
\teil Ein zylindrischer Regenwassertank fasst $V = 50\,\text{m}^3$ bei einer Höhe von
$h = 2{,}5\,\text{m}$. Berechne den Radius auf zwei Stellen genau. Es gilt
$V = \pi \cdot r^2 \cdot h$. \hfill (P10 2022)
\end{teile}
\end{aufgabe}

\begin{aufgabe}{Wurzelziehen -- Aussagen prüfen. Kreuze für jede Aussage an, ob sie
wahr oder falsch ist. \hfill (P10 2021)}

\sachtabelle{lcc}{Aussage & wahr & falsch}{%
a) \ $(x - 5)^2 = 0$ hat genau eine Lösung. & $\square$ & $\square$ \\
b) \ $5$ und $13$ sind Lösungen von $(x + 9)^2 = 16$. & $\square$ & $\square$}

\begin{teile}
\setcounter{teil}{2}
\teil Gib eine quadratische Gleichung an, die keine Lösung hat. \hfill \leerfeld
\end{teile}
\end{aufgabe}
